\documentclass[11pt]{article}
\usepackage[a4paper,margin=25mm]{geometry}
\usepackage[T1]{fontenc}
\usepackage[utf8]{inputenc}
\usepackage[provide=*,czech]{babel}
\usepackage{lmodern,amsmath,amssymb,amsthm,booktabs,array,needspace,hyperref}
\setlength{\emergencystretch}{2em}
\hypersetup{colorlinks=true,linkcolor=blue,urlcolor=blue,
pdftitle={Asymptotické otevřené problémy o férových kostkách},
pdfauthor={Výzkumné poznámky k rukopisu Fair-Dice-writeup}}
\newtheorem{veta}{Věta}
\newcommand{\soubor}[1]{\par\smallskip\noindent\path{#1}\par\smallskip}
\title{Asymptotické otevřené problémy o férových kostkách\\
\large Výsledky a mapa důkazů}
\author{Výzkumné poznámky k rukopisu \texttt{Fair-Dice-writeup}}
\date{30. září--1. října 2026}
\begin{document}\maketitle

Nejvýraznější výsledek se týká třetí otázky v závěru článku:
\emph{exponenciální dolní mez pro každou jednotlivou kostku neplatí.}
Navíc lze určit správný asymptotický řád minima pro jednu zadanou kostku:
je to $\Theta(n^2)$. Otázky optimální celkové velikosti a polynomiálního
vyrovnání počtů stěn zůstávají otevřené.

Zpráva se soustředí na obecná tvrzení pro rostoucí $n$. Hlavní důkazy o kvadratickém minimu, více malých kostkách a koncovém profilu prošly
třemi oddělenými matematickými kontrolami spolupracujících agentů. Jde o interní kontrolu důkazů,
nikoli strojovou formalizaci nebo externí recenzní řízení. Rešerše
neopravňuje k definitivnímu tvrzení o historické novosti všech výsledků.
Zdrojový článek nebyl v rámci této práce přepisován.

\section{Tři původní otázky}

Označme $A_n$ nejmenší celkový počet stěn v permutačně férové rodině
$n$ konečných rovnoměrných kostek a $E_n$ stejné minimum při požadavku
stejných počtů stěn. Dále nechť $g(n)$ je nejmenší počet stěn jedné
zadané kostky; ostatní kostky mohou být libovolně velké.

\begin{center}
\begin{tabular}{>{\raggedright\arraybackslash}p{.31\textwidth}>{\raggedright\arraybackslash}p{.61\textwidth}}
\toprule
Otázka & Výsledek\\
\midrule
Asymptotika $A_n$ & Stále $\exp(\Omega(n))\le A_n\le
\exp(O(n\log n))$. Mezera nebyla uzavřena.\\[4pt]
Je $E_n\le A_n^C$? & Otevřeno. Dostáváme alespoň
$E_n\le A_n^{O(\log n)}=A_n^{O(\log\log A_n)}$.\\[4pt]
Musí být každá kostka exponenciální? & Ne. Platí
$g(n)=\Theta(n^2)$, a dokonce lze předepsat každý dostatečně velký
kvadratický počet stěn.\\
\bottomrule
\end{tabular}
\end{center}

Uvedený kvazipolynomiální odhad pro $E_n$ je elementárním důsledkem
již známé absolutní horní meze $\exp(O(n\log n))$ a exponenciální
dolní meze pro $A_n$. Nejde o nový vyrovnávací postup aplikovaný na
optimální nevyrovnanou rodinu. Konstrukce se $C(n!)^2$ stěnami na každé
kostce zpřesňuje explicitní horní mez, ale nemění její známý asymptotický řád.
\soubor{size_bounds/factorial_squared_construction.tex}
\soubor{size_bounds/equalization_notes.tex}

\Needspace{11\baselineskip}
\section{Jedna zadaná kostka: přesně kvadratický řád}

\begin{veta}
Existují absolutní konstanty $c,C>0$ takové, že:
\begin{enumerate}
\item každá jednotlivá kostka v permutačně férové rodině $n$ konečných
rovnoměrných kostek má alespoň $cn^2$ stěn;
\item pro každé celé $K\ge Cn^2$ existuje taková rodina, v níž má
zadaná kostka přesně $K$ stěn;
\item ostatních $n-1$ kostek v této konstrukci může mít jeden společný
konečný počet stěn.
\end{enumerate}
\end{veta}

Počet stěn ostatních kostek v horní konstrukci může být velmi velký.
Věta tedy neřeší minimum součtu stěn a je slučitelná s tím, že téměř
všechny kostky musejí být exponenciálně velké.

\subsection*{Myšlenka horní konstrukce}

Reálná rovnoměrná kvadratura stupně $O(n)$ potřebuje jen $O(n^2)$ uzlů.
Existenci pro každý dostatečně velký počet uzlů dává rovnoměrný případ
věty Gilboy a Peleda; viz
\href{https://arxiv.org/abs/1507.01505}{Chebyshev-type quadratures for doubling weights,
Theorem 1.1}. Uzly se nejprve regularizují, aby byly různé a uvnitř
intervalu. Poté se aproximují racionálními čísly.

Vzniklou malou chybu neopravujeme přidáváním stěn zadané kostce.
Opravují ji malé racionální změny hustot ostatních kostek, lokalizované
kolem různých uzlů. Každá změna má nulové plné momenty do potřebného
stupně. Díky disjunktním podporám v integrálu pořadí zaniknou všechny
vyšší smíšené členy. Zbývající lineární člen je přesně derivace chyby
kvadratury. Vandermondovy soustavy umožní tuto chybu odstranit.

Nakonec jsou všechny hustoty kladné racionální schodové funkce.
Každý interval konstantních hustot lze nahradit již známým konečným
férovým blokem s vhodně násobenými počty jednotlivých písmen.
Tím vznikne rodina skutečných konečných rovnoměrných kostek.

\soubor{size_bounds/polynomial_exceptional_die.tex}
\soubor{size_bounds/centered_bump_simplification.tex}
\soubor{rational_designs/asymptotic/every_large_exceptional_size.tex}

\subsection*{Myšlenka dolní meze}

Uspořádáme stěny zadané kostky a zapíšeme $p_i(q)$ pro podíl stěn
$i$té jiné kostky pod její $q$tou stěnou. Sloupce jsou neklesající a
férovost dává přesné smíšené momenty
\[
 \mathbb E_w\prod_{i\in S}p_i=\frac1{|S|+1},\qquad m=n-1.
\]
Pro tuto část důkazu dokonce stačí, aby zadaná kostka vyhrávala každý
podsoubor, který ji obsahuje, s pravděpodobností převrácené velikosti
tohoto podsouboru.

Závěrečná kontrola přinesla krátký elementární důkaz. Monotónní řádky
spojíme pomocnou spojitou křivkou. Pro dva sloupce a bod této křivky
je součin rozdílů od jejich hodnot v tomto bodě nezáporný na všech
řádcích. Sloupce rozdělíme do disjunktních dvojic; v každé dvojici
zvolíme první okamžik, kdy některý sloupec dosáhne předepsané hodnoty.
Jeden kořen příslušného skalárního součinu je pak přesně předepsaný,
druhý leží pod ním. Přesné smíšené momenty dovolují jeho střední
hodnotu vypočítat jednorozměrným integrálem.

Gaussova kvadratura nejprve ukáže, že v posledním řádku je
každé $p_i$ alespoň $1-C/m^2$. Gaussova--Radauova kvadratura
s koncovým uzlem $1$ pak ze součinu dvojic vytvoří normalizovaný nezáporný test,
jehož integrál je $O(m^{-2})$ a jehož hodnota na posledním řádku
je zdola omezená. Proto je váha posledního řádku $O(m^{-2})$.
Jelikož u rovnoměrné kostky je váha každé stěny $1/K$,
dostáváme $K\ge cm^2$. Důkaz dává například explicitní mez
$K\ge(n-1)^2/200$.

Potřebné kvadraturní skutečnosti jsou v důkazu odvozeny z elementárních
identit Legendreových polynomů. Starší analytický důkaz pomocí
Laguerrových polynomů zůstává uložený; pro tento výsledek už není nutný.
\soubor{size_bounds/QUADRATIC_PROOF_GUIDE.txt}
\soubor{independent_directions/elementary_quadratic_atom_bound.tex}
\soubor{size_bounds/quadratic_per_die_lower_bound.tex}

\section{Více zadaných malých kostek}

\begin{veta}
Pro každé pevné $k$ existují konstanty $C_k,c_k>0$ takové, že pro
všechna dostatečně velká $n$ má některá úplně konečná permutačně
férová rodina $k$ zadaných kostek se stejným počtem stěn, nejvýše
$C_k n^{c_k}$. Ostatních $n-k$ kostek může mít jiný společný
konečný počet stěn.
\end{veta}

Závislost na $k$ se podařilo řídit jednotně. Existuje absolutní $C$
takové, že pro $n>k\ge1$ lze společný počet stěn zadaných kostek
omezit výrazem
\[
 S\le(n+2^k+2)^{\exp(Ck^2)}.
\]
Pro nějaké absolutní $c>0$ tedy může mít
$\lfloor c\sqrt{\log n}\rfloor$ zadaných kostek každá nejvýše
$\exp(\sqrt n)$ stěn. Pro $k=o(\sqrt{\log n})$ dostáváme
$S\le\exp(n^{o(1)})$. Jde o subexponenciální velikosti rostoucího
počtu zadaných kostek; ostatní kostky mohou zůstat velmi velké.
Jednotný polynomiální odhad $n^C$ pro rostoucí $k$ tím získán není.
Pro dvě kostky dává samostatná,
kvantitativnější konstrukce počty $O(n^6)$ a $O(n^{162})$; tyto
exponenty rovněž nebyly optimalizovány.

Důkaz postupně odděluje z pevného počtu původně rovnoměrných
spojitých proměnných jednotlivé hustoty. Jejich podpory leží v
čistých intervalech: v každém intervalu je kladná právě jedna
z těchto hustot. Přidávaná proměnná proto může ve stejném intervalu
potkat nejvýše jednu starší proměnnou. Pro zachování všech pořadí
stačí lokálně zachovat uspořádané polynomiální momenty dvojice.

Klíčové je společně řídit jmenovatele napříč všemi intervaly.
Při každém kroku se udržuje zásoba nízkých momentů na jednom
polynomiálním racionálním rastru. Pro získání výstupních momentů do
stupně $q$ stačí vstupní zásoba do stupně $2q+1$. Konečný počet kroků
proto vyžaduje jen pevnou počáteční zásobu. Oddělené Legendreovy
perturbace zajišťují polynomiální podmíněnost momentových soustav.
Stejný malý racionální podblok soustavy ve všech intervalech pak
zajišťuje jeden společný jmenovatel; samotné oddělené odhady
jmenovatelů jednotlivých řešení by nestačily.

Kladnost při dokončování hustot se dokazuje pomocí podílového odhadu
polynomů modulo nízké momenty v jednotlivých intervalech. Všechny
intervaly lze udržet polynomiálně široké. Kvadratura uvnitř čistých
intervalů a zobecněná lokální korekce nakonec převedou celou rodinu
na skutečné konečné kostky, při zachování polynomiálních počtů stěn
zadaných $k$ kostek.

\soubor{independent_directions/fixed_number_polynomial_dice.tex}
\soubor{size_bounds/synchronized_local_moment_completion.tex}
\soubor{size_bounds/uniform_grid_moment_refinement.tex}
\soubor{size_bounds/clustered_multi_exceptional_lifting.tex}
\soubor{size_bounds/two_polynomial_exceptional_dice.tex}
\soubor{independent_directions/FIXED_K_PROOF_GUIDE.txt}
\soubor{size_bounds/uniform_prefix_conditioning.tex}
\soubor{size_bounds/uniform_fixed_number_complexity.tex}

\section{Další asymptotické výsledky}

\paragraph{Všechny dostatečně velké přípustné velikosti.}
Položme $P_n=\prod_{p\le n}p$. Existuje absolutní konstanta $C$ taková,
že každý celý počet $M\ge\exp(Cn\log(n+1))$ dělitelný $P_n$ lze
realizovat jako společný počet stěn permutačně férových $n$ kostek.
Nad touto hranicí tedy již nejsou další aritmetické překážky.

Pomocná konstrukce doplní každé počáteční slovo $u$ doslova jako prefix
férového slova se společným počtem stěn nejvýše
$(|u|+1)\exp(O(n\log(n+1)))$.
Přesnost podpisu slova se při jednotlivých krocích zhruba zdvojnásobuje.
Relabelování vychází z publikované věty Kuperberga, Lovetta a Peleda o
\href{https://arxiv.org/abs/1302.4295}{$t$-rovnoměrných množinách permutací};
stejné součty mocnin řídí pozice bloků a zruší přepravované chyby.
Krátký podepsaný férový cíl a růstový odhad pro kladné zlomky pak dají
dvě realizovatelné velikosti $A,A+P_n$. Jejich kladné kombinace vyplní
všechny dostatečně velké násobky $P_n$. Minimum celkové velikosti tím
zůstává mezi původními dvěma asymptotickými řády.
\soubor{size_bounds/order_doubling_prefix_completion.tex}
\soubor{size_bounds/factorial_order_eventual_conductor.tex}

\paragraph{Co brání jednoduchému rozměrovému důkazu.}
Pro férová slova, jejichž parametrizace vahami stěn má plnou hodnost
v prostoru redukovaných Lieových podpisů, je optimální délka
$\exp(\Theta(n\log n))$. Případná konstrukce délky $\exp(O(n))$
tedy musí mít singulární parametrizaci. Tento odhad se na všechna
férová slova nevztahuje. Obecné konstrukce ukazují, že i velmi mnoho
různých uspořádaných prefixových souřadnic může být totožných.
Samotné momenty na disjunktních množinách písmen navíc připouštějí
kladné semidefinitní doplnění s malou hodností. Silnější dolní mez
pro celkovou velikost by musela využít další strukturu skutečného slova.
\soubor{size_bounds/regular_fair_realizations_and_rank_barrier.tex}
\soubor{independent_directions/prefix_signature_rank_obstruction.tex}
\soubor{rational_designs/asymptotic/squarefree_prefix_gram_barrier.tex}
\soubor{size_bounds/fair_signature_hankel_rank.tex}

\paragraph{Přibližná férovost.}
Jestliže pro každé úplné pořadí $\pi$ platí
$|n!\Pr(\pi)-1|\le\varepsilon\le1/2$, pak má každá kostka alespoň
\[
 c n\min\{n,\log(1/\varepsilon)\}
\]
stěn. Pro pevnou chybu je tato mez lineární; pro exponenciálně malou
chybu je kvadratická. Důkaz přesné dolní meze lze stabilizovat, protože
chyba polynomiálního testu stupně $d$ roste jako
$\varepsilon\exp(O(d))$, bez dodatečného faktoru $n^d$.
Při $\varepsilon=0$ rozumíme $\log(1/\varepsilon)=+\infty$.
\soubor{rational_designs/asymptotic/sharp_approximate_per_die_lower_bound.tex}

Randomizovaná konstrukce z palindromických bloků dává opačný odhad
\[
 O\!\left(n^{13/10}(\log n)^{1/5}\varepsilon^{-2/5}\right)
\]
stěn na kostku. Pro pevné $\varepsilon$ tedy zůstává zajímavá mezera
mezi lineární dolní mezí a touto horní mezí.
\soubor{approximate_random/refined_random_palindrome_bound.tex}

\paragraph{Racionální skalární kvadratura.}
Je-li $R_r$ nejmenší počet stejně vážených racionálních uzlů kvadratury
rovnoměrné míry na intervalu, přesné do stupně $r$, přičemž uzly se
mohou opakovat, platí
\[
 2^{2\lfloor r/2\rfloor}\le R_r\le2^r+O(2^{r/2}).
\]
Zejména pro sudé $r$ je $R_r/2^r\to1$. Horní mez používá publikovanou
lokální teorii a kruhovou metodu; nejde o tvrzení o novosti této
číselněteoretické části. Tento výsledek vysvětluje, proč protipříklad
k exponenciálnímu počtu stěn jednotlivé kostky musel využít rozdílné
srovnávací sloupce, nikoli malou racionální skalární kvadraturu.
\soubor{independent_directions/optimized_two_adic_filter.tex}
\soubor{size_bounds/mitkin_sharp_rational_upper.tex}

\paragraph{Struktura srovnávacích sloupců.}
V obecném monotónním modelu smíšených momentů položme
$u_i=1-p_i$ a $s=\sum_i u_i$. Platí ostrý relativní odhad
\[
 \max_i m(1-p_i)\le C(s+1/m),\qquad s=\sum_i(1-p_i).
\]
Navíc pro každé pevné $L<\infty$ platí jednotně až ke kraji
\[
 \max_i|m(1-p_i)-s|\le C_L R_m(s),\qquad
 R_m(s)=\min\{ms,\sqrt{s/m}+1/m\},
 \qquad 0\le s\le L.
\]
Na každém pevném kladném rozsahu je tedy chyba $O(m^{-1/2})$.
Pro každou škálu $\lambda\in(0,L]$ existují skutečné konečné rovnoměrné
férové kostky s řádkem, kde $s\asymp_L\lambda$ a chyba je alespoň
$c_L R_m(\lambda)$. Obálka je tedy ostrá na všech omezených koncových
škálách. Příklady předepisují $s$ až na konstantní násobek, nikoli jeho
přesnou hodnotu. Gaussovy uzly dávají větev $\lambda\ge1/m$;
Legendreova perturbace s proměnnou amplitudou pokrývá menší škály.
Přenos na skutečné kostky neposkytuje odhad počtu stěn: aproximuje
skalární momentovou míru pravidelnou kvadraturou a využívá libovolně
malých korekcí v konstrukci jedné zadané kostky.
Z horního odhadu také plyne relativní sbližování
$m(1-p_i)/s\to1$, stejnoměrně ve všech sloupcích, kdykoli $ms\to\infty$
a $s\le L$.
Pro obecné kladné váhy řádků navíc platí
$w_q\le C_L(\sqrt{s_q}/m^{3/2}+m^{-2})$.
Ostrost této poslední meze je doložena ve váženém modelu;
odpovídající optimální počty stěn rovnoměrných kostek z ní nevyvozujeme.
\soubor{sharp_endpoint_relative_balance.tex}
\soubor{endpoint_profile_sharpness.tex}
\soubor{endpoint_profile_fluctuation_barrier.tex}
\soubor{finite_profile_realization.tex}
\soubor{rational_designs/asymptotic/quantitative_uniform_endpoint_profile.tex}
\soubor{independent_directions/optimal_interior_endpoint_profile.tex}
\soubor{rational_designs/asymptotic/variable_jackson_endpoint_kernel.tex}
\soubor{independent_directions/natural_scale_endpoint_profile.tex}
\soubor{independent_directions/intermediate_endpoint_sharpness.tex}
\soubor{rational_designs/asymptotic/natural_scale_row_atoms.tex}
\soubor{independent_directions/ENDPOINT_PROFILE_PROOF_GUIDE.txt}

\section{Otázky, které nyní vypadají nejzajímavěji}

\begin{enumerate}
\item \textbf{Původní mezera pro celkovou velikost.}
Lze dosáhnout $\exp(O(n))$ celkových stěn, nebo platí silnější
dolní mez? Dosavadní aritmetické překážky samy rozlišení nedávají.

\item \textbf{Polynomiální vyrovnání.}
Lze z nerovnoměrných počtů stěn přejít ke stejným počtům za
polynomiální cenu? Samotné vyrovnání četností bloků nestačí:
pořadí výběrů z různých bloků může férovost porušit.

\item \textbf{Ostré minimum pro dvě zadané kostky.}
Mohou mít obě jen $O(n^2)$ stěn? Současná konstrukce dává
polynomiální, ale velmi hrubou horní mez. Metoda s alespoň $m^2$
uzly v každé čisté skupině má vlastní kubickou bariéru; její prosté
optimalizování proto ke kvadratickému výsledku nestačí.
Viz \path{size_bounds/clustered_lifting_cubic_barrier.tex}.

\item \textbf{Rostoucí počet malých kostek.}
Pro každé pevné $k$ je polynomiální existence vyřešena a exponent lze
omezit $\exp(O(k^2))$. Lze tuto závislost výrazně zlepšit?
Jsou-li všechny malé kostky omezeny
stejným $n^C$, pak nutně $k\le(2C+o(1))\log n$.
Pokud by šlo dosáhnout $c\log N$ kostek s nejvýše $N^C$ stěnami,
s absolutními $c,C>0$, omezování na jejich podsoubory by už dalo
$A_r=\exp(O(r))$ a vyřešilo první původní otázku. Současná věta
pro pevné $k$ takový závěr nedává. Viz
\path{many_small_dice_reduction.tex}.

\item \textbf{Pevná přibližná chyba.}
Stačí $n^{1+o(1)}$ stěn na každé kostce? Nebo je univerzálně nutný
větší než lineární řád? Výsledky pro jednu konkrétní randomizovanou
konstrukci neposkytují dolní mez pro všechny konstrukce.

\item \textbf{Profil na rostoucím rozsahu.}
Na každém pevném omezeném rozsahu $s$ je již dokázán odhad
$C_L(\sqrt{s/m}+1/m)$, včetně kraje $s=0$.
Jak optimálně závisí konstanta $C_L$ na $L$? Lze obdobně ostré
sbližování popsat pro rostoucí $s=s(m)$ a postupně až ve středu profilu?
\end{enumerate}

\section{Jak číst uložené podklady}

Pro první čtení doporučuji horní konstrukci jedné malé kostky,
poté \path{size_bounds/QUADRATIC_PROOF_GUIDE.txt} a závěrečnou kvadratickou
dolní mez. Důkazy pro pevný počet malých kostek tvoří samostatnou větev a nejsou
závislostí výsledku $g(n)=\Theta(n^2)$.
Jejich nejkratší mapa je
\path{independent_directions/FIXED_K_PROOF_GUIDE.txt}.
Jednotnou závislost na $k$ rozepisuje také
\path{size_bounds/FIXED_K_PROOF_GUIDE.txt}.

Soubory \path{ASYMPTOTIC_STATUS.txt}, \path{PROGRESS.txt} a
\path{ROOT_REVIEW.txt} zachycují aktuální stav, průběh a kontroly.
Stručný audit výsledného balíku je v \path{FINAL_REVIEW.txt}.
Ve složce jsou zachovány i starší mezivýsledky a neúspěšné směry.
Jejich přítomnost neznamená, že jsou dalšími potřebnými částmi
aktuálního důkazu. U výsledků v této zprávě existují souběžné zdrojové
soubory \texttt{.tex} a zkompilované \texttt{.pdf}.

\end{document}
